\chapterhead{58}{ORR-19}{Range of Practices and Issues in Economic Capital Frameworks}{4}{1}

\begin{ornotebox}{About this chapter}
The crash course original covers only risk measures and risk aggregation. The
remainder of LO 58.a, and all of LO 58.b, c and d, are re-typeset here from the
class lecture notes that complete this reading.
\end{ornotebox}

\lohead{58}{a}{Within the economic capital implementation framework, describe the
challenges that appear in: defining and calculating risk measures; risk
aggregation; validation of models; dependency modeling in credit risk; evaluating
counterparty credit risk; and assessing interest rate risk in the banking book.}

\orsub{Risk measures and risk aggregation}

\orsubb{Key properties of a coherent risk measure}

\begin{lettered}
\item \textbf{Monotonicity:} higher returns suggest lower risk.
\item \textbf{Subadditivity:} risk of a portfolio $\leq$ sum of individual asset
risks.
\item \textbf{Positive homogeneity:} risk scales with the portfolio size.
\item \textbf{Translation invariance:} risk depends on portfolio contents.
\end{lettered}

\orsubb{Ideal characteristics of risk measures}

\begin{center}
\begin{tikzpicture}[font=\fontsize{8.2}{10}\selectfont,
  ch/.style={text width=4.6cm,align=center,rounded corners=3pt,inner sep=6pt,
             fill=paleblue,draw=palenavy,line width=0.6pt,text=navy}]
\node[ch] at (-5.2, 0.55) {a) Intuitive};
\node[ch] at ( 0.0, 0.55) {b) Stable};
\node[ch] at ( 5.2, 0.55) {c) Easy to compute and understand};
\node[ch] at (-5.2,-0.85) {d) Coherent};
\node[ch] at ( 0.0,-0.85) {e) Economically interpretable};
\node[ch] at ( 5.2,-0.85) {f) Allows simple risk decomposition};
\end{tikzpicture}
\end{center}
\orfigcap{Six characteristics of an ideal risk measure.}

\orsubb{Common risk measures}

\noindent\begin{minipage}{\textwidth}
\renewcommand{\arraystretch}{1.35}
\begin{tabularx}{\textwidth}{@{}L{0.28\textwidth}Y@{}}
\rowcolor{navy} \thd{Measure} & \thd{Assessment}\\
{\tabfont \textbf{1. Standard deviation}} & {\tabfont Easy to calculate but not coherent; violates monotonicity.}\\
\rowcolor{paleblue}
{\tabfont \textbf{2. Value at Risk (VaR)}} & {\tabfont Most commonly used but not coherent; violates subadditivity.}\\
{\tabfont \textbf{3. Expected Shortfall (ES)}} & {\tabfont More stable than VaR but less interpretable.}\\
\rowcolor{paleblue}
{\tabfont \textbf{4. Spectral / distorted measures}} & {\tabfont Not intuitive or commonly used.}\\
\end{tabularx}
\end{minipage}
\orfigcap{Four risk measures, and where each falls short.}

\orsubb{Challenges of risk measures}
\begin{lettered}
\item No single measure perfectly captures all aspects of risk.
\item VaR and ES are the most common measures, but neither is perfect.
\end{lettered}

\orsubb{Risk aggregation}

Involves combining various risk types (market, credit, operational, business) into
a single measure. Challenges include:

\begin{lettered}
\item Classifying risks into types (market, credit, operational, business).
\item Aggregating risks across different time horizons and confidence levels.
\item Accounting for dependencies between different risk types.
\end{lettered}

\orsub{Validation of models}

\begin{orkeybox}
Validation is the process of proving that a model works as intended. It is good
for testing risk sensitivity, but less effective for checking extreme losses (tail
events, high quantiles). Economic capital models are harder to validate because
they produce a loss \emph{distribution}, not a single forecast.
\end{orkeybox}

\orsubb{Qualitative validation}

\noindent\begin{minipage}{\textwidth}
\renewcommand{\arraystretch}{1.3}
\begin{tabularx}{\textwidth}{@{}L{0.24\textwidth}YY@{}}
\rowcolor{navy} \thd{Process} & \thd{Purpose} & \thd{Main challenge}\\
{\tabfont \textbf{1. Use test}} & {\tabfont Regulators rely more on models if the bank uses them for internal decisions.} & {\tabfont Regulators struggle to see which model properties are actually used versus ignored.}\\
\rowcolor{paleblue}
{\tabfont \textbf{2. Qualitative review}} & {\tabfont Review documentation, methodology, algorithms, and risk drivers.} & {\tabfont Ensuring the model works in theory, captures correct risk drivers, and has accurate maths.}\\
{\tabfont \textbf{3. Systems implementation}} & {\tabfont Conduct User Acceptance Testing (UAT) and code checks before rolling out.} & {\tabfont Ensure no implementation errors.}\\
\rowcolor{paleblue}
{\tabfont \textbf{4. Management oversight}} & {\tabfont Senior management reviews and uses model outputs.} & {\tabfont Management must correctly interpret results.}\\
{\tabfont \textbf{5. Data quality checks}} & {\tabfont Ensure data is accurate, complete, and relevant.} & {\tabfont Detect missing or incorrect data.}\\
\rowcolor{paleblue}
{\tabfont \textbf{6. Sensitivity testing}} & {\tabfont Test assumptions (correlations, recovery rates, tail distributions).} & {\tabfont Assess impact of assumption changes.}\\
\end{tabularx}
\end{minipage}
\orfigcap{Six qualitative validation processes.}

\orsubb{Quantitative validation processes}

\noindent\begin{minipage}{\textwidth}
\renewcommand{\arraystretch}{1.3}
\begin{tabularx}{\textwidth}{@{}L{0.26\textwidth}YY@{}}
\rowcolor{navy} \thd{Process} & \thd{Purpose} & \thd{Main challenge}\\
{\tabfont \textbf{1. Validate inputs and parameters}} & {\tabfont Check model inputs such as correlations.} & {\tabfont Correct inputs alone cannot eliminate model error.}\\
\rowcolor{paleblue}
{\tabfont \textbf{2. Model replication}} & {\tabfont Reproduce model results independently.} & {\tabfont Simply re-running the same algorithms is insufficient; rarely used in practice.}\\
{\tabfont \textbf{3. Benchmarking}} & {\tabfont Compare with standard models or reference portfolios.} & {\tabfont Shows consistency, not necessarily accuracy.}\\
\rowcolor{paleblue}
{\tabfont \textbf{4. Backtesting}} & {\tabfont Compare predicted outcomes with actual outcomes.} & {\tabfont Difficult for economic capital models with long horizons.}\\
{\tabfont \textbf{5. P\&L attribution}} & {\tabfont Compare actual profits and losses with model risk drivers.} & {\tabfont Not widely used outside of market risk pricing models.}\\
\rowcolor{paleblue}
{\tabfont \textbf{6. Stress testing}} & {\tabfont Test model under extreme scenarios.} & {\tabfont Focuses on extreme scenarios rather than day-to-day absolute accuracy.}\\
\end{tabularx}
\end{minipage}
\orfigcap{Six quantitative validation processes.}

\orsub{Dependency modeling in credit risk}

Measures how defaults of different borrowers are related. Important because
defaults are not independent, especially during economic downturns. In general,
dependencies can be modelled using credit risk portfolio models, models using
copulas, and models based on the asymptotic single-risk factor (ASRF) model.

\orsubb{Asymptotic Single-Risk Factor (ASRF) model}

The foundation for the IRB framework under Basel II to compute capital
requirements.

\begin{checks}
\item Assumes a single systematic risk factor and infinite portfolio granularity.
\item \textbf{Challenge:} banks must justify how they handle the single-factor
limitation and how they calibrate correlations.
\end{checks}

\orsubb{Model assumptions}

Many assumptions may not hold in reality:

\begin{enumerate}
\item[\textcolor{rust}{\sffamily\bfseries 1)}] Gaussian copula (normal
dependence).
\item[\textcolor{rust}{\sffamily\bfseries 2)}] Normally distributed default
drivers.
\item[\textcolor{rust}{\sffamily\bfseries 3)}] Stable correlations over time.
\item[\textcolor{rust}{\sffamily\bfseries 4)}] Common risk factors fully explain
default correlation.
\end{enumerate}

\orsubb{Major model weaknesses and impacts: the tail-risk and clustering problem}

\begin{lettered}
\item \textbf{Underestimation of capital.} Traditional models fail to explain the
time-clustering of defaults seen in real markets.
\item \textbf{PD/LGD correlation.} Models often fail to integrate the correlation
between Probability of Default and Loss Given Default. Ignoring LGD variability
leads to severe underestimation of required economic capital.
\end{lettered}

\orsubb{Data source errors}

\begin{itemize}
\item \textbf{Business cycle bias.} Rating changes vary wildly between expansions
and recessions. If the sample calibration period is poorly chosen, correlation
estimates will be heavily distorted.
\item \textbf{The equity price proxy problem.} Many models use observable equity
prices to approximate unobservable asset returns.
  \begin{itemize}
  \item \emph{The flaw:} this relationship is unobservable and often nonlinear.
  \item \emph{The impact:} equity prices contain noise irrelevant to credit risk,
  causing inaccurate correlation estimates.
  \end{itemize}
\end{itemize}

\orsubb{Regulatory approach (Basel II) challenges}

When banks rely strictly on regulatory-type models, they face unique hurdles:

\begin{checks}
\item \textbf{Data scarcity:} limited historical data makes estimating
correlations difficult.
\item \textbf{Model inconsistency:} the bank's internal assumptions often conflict
with the underlying maths of the Basel II model.
\end{checks}

\orsub{Evaluating counterparty credit risk}

Such a task is a significant challenge because it requires obtaining data from
multiple systems, measuring exposures from an enormous number of transactions
(including many that exhibit optionality) spanning a wide range of time periods,
monitoring collateral and netting arrangements, and categorising exposures across
many counterparties.

\orsubb{Market risk-related challenges to counterparty Exposure at Default (EAD) estimation}

Counterparty exposure uses market risk simulations, similar to VaR, but with two
key differences:

\begin{enumerate}
\item[\textcolor{rust}{\sffamily\bfseries 1)}] No netting across counterparties is
allowed; exposures must be calculated at the netting set level, increasing
computational complexity.
\item[\textcolor{rust}{\sffamily\bfseries 2)}] Unlike VaR's single short-term
horizon, EAD requires multiple future time horizons, meaning market factors and
exposures must be simulated far into the future.
\end{enumerate}

\orsubb{Credit risk-related challenges to PD and LGD estimation}

Banks must estimate Probability of Default (PD) and Loss Given Default (LGD) for
each counterparty and transaction. This is particularly difficult for hedge funds,
where information on leverage, volatility, and investment strategy is often
limited. Even for counterparties with existing relationships, banks still need
transaction-specific LGD estimates.

\orsubb{Interaction between market risk and credit risk: wrong-way risk}

\begin{ordefbox}{Wrong-way risk}
Wrong-way risk occurs when exposure increases as the counterparty's credit quality
deteriorates.
\end{ordefbox}

\begin{itemize}
\item Measuring it requires understanding the market risk factors affecting the
counterparty and comparing them with the bank's own exposures.
\item It is difficult to identify and quantify, especially for less transparent
counterparties, and requires long-horizon, high-confidence economic capital models.
\end{itemize}

\orsubb{Operational risk-related challenges in managing counterparty credit risk}

\begin{itemize}
\item Managing counterparty credit risk requires specialised systems and
experienced personnel.
\item Complex activities such as daily limit monitoring, mark-to-market, collateral
management, and intraday liquidity and credit monitoring increase the likelihood
of operational errors.
\item Quantifying operational risk is particularly difficult for new products,
rapidly growing businesses, and rare but severe events.
\end{itemize}

\orsubb{Differences in risk profiles between margined and non-margined counterparties}

The main difference lies in the forecasting period:

\begin{lettered}
\item \textbf{Margined counterparties:} short forecasting horizon due to frequent
collateral exchanges.
\item \textbf{Non-margined counterparties:} longer forecasting horizon because
exposures remain uncollateralised for longer.
\end{lettered}

Using different time horizons makes risk aggregation more difficult.

\orsubb{Aggregation challenges}

Moving from single-counterparty risk to firm-wide economic capital significantly
increases complexity. Aggregating exposures across derivatives, securities
financing, market risk, credit risk, and operational risk requires enormous
computational power and data, and many banks still struggle due to outdated
systems.

\orsub{Assessing interest rate risk in the banking book}

\orsubb{Optionality in the banking book}

A major challenge is measuring nonlinear risk from embedded options.

\begin{lettered}
\item \textbf{Asset side:} products such as mortgages, mortgage-backed securities
(MBS), and consumer loans contain prepayment options, creating uncertain cash
flows and making interest rate risk difficult to measure.
\item \textbf{Liability side:} non-maturity deposits contain two embedded options:
  \begin{itemize}
  \item The bank can decide the interest rate paid and when to change it.
  \item The depositor can withdraw funds at any time without penalty.
  \end{itemize}
\end{lettered}

The interaction of these options creates significant valuation and interest rate
sensitivity challenges, requiring complex stochastic-path models.

\orsubb{Banks' pricing behaviour}

Models must estimate the persistence of non-maturity deposits and determine deposit
rates based on market conditions, customer relationships, competitive power, and
business policies. Pricing must also incorporate credit risk, requiring models that
link credit spreads to changes in macroeconomic conditions and interest rates.
Therefore, interest rate stress scenarios should jointly consider interest rate and
credit risk.

\orsubb{Choice of stress scenarios}

Using simple interest rate shocks has several limitations:

\begin{checks}
\item They are not probability-based, making them difficult to integrate into
VaR-based economic capital models.
\item They may not reflect the current interest rate or economic environment.
\item They ignore changes in the slope and curvature of the yield curve.
\item They do not allow an integrated analysis of interest rate risk and credit
risk in the banking book.
\end{checks}

% ---------------------------------------------------------------------
\lohead{58}{b}{Describe the recommendations by the Bank for International
Settlements (BIS) that supervisors should consider in order to make effective use
of internal risk measures, such as economic capital, that are not designed for
regulatory purposes.}

\orsub{BIS recommendations for supervisors}

\begin{enumerate}
\item[\textcolor{rust}{\sffamily\bfseries 1.}] \textbf{Use in capital adequacy.}
Economic capital models should be actively used in decision-making to assess
capital adequacy. The board should understand the difference between gross
(stand-alone) risk and net (diversified) risk while setting the bank's risk
tolerance.
\item[\textcolor{rust}{\sffamily\bfseries 2.}] \textbf{Senior management.} Senior
management should provide strong commitment and oversight. They must understand the
importance of economic capital and ensure adequate governance, infrastructure, and
resources are in place.
\item[\textcolor{rust}{\sffamily\bfseries 3.}] \textbf{Transparency and
decision-making.} Economic capital results should be clear, transparent, and easy
to interpret. The models should support business decisions and provide flexibility
for firm-wide stress testing.
\item[\textcolor{rust}{\sffamily\bfseries 4.}] \textbf{Risk identification.} Banks
should identify all relevant risk drivers, positions, and exposures to ensure
accurate measurement of economic capital. Risks that are difficult to quantify
should be assessed using stress testing, scenario analysis, or sensitivity
analysis.
\item[\textcolor{rust}{\sffamily\bfseries 5.}] \textbf{Risk measures.} There is no
universally best risk measure. Banks should understand the strengths and
limitations of their chosen risk measures before using them for economic capital.
\item[\textcolor{rust}{\sffamily\bfseries 6.}] \textbf{Risk aggregation.} Risk
aggregation should combine different risk types using consistent methodologies. The
aggregation process should reflect the bank's business model, risk profile, and
interrelationships among risks.
\item[\textcolor{rust}{\sffamily\bfseries 7.}] \textbf{Validation.} Economic
capital models should be thoroughly validated through multiple tests to ensure they
perform as intended. The estimated capital should be sufficient to absorb
unexpected losses at the selected confidence level.
\item[\textcolor{rust}{\sffamily\bfseries 8.}] \textbf{Dependency modeling.} Banks
should assess whether the credit dependency (correlation) assumptions used in their
models are appropriate. Model limitations should be addressed through stress
testing and scenario analysis.
\item[\textcolor{rust}{\sffamily\bfseries 9.}] \textbf{Counterparty credit risk.}
Banks should select an appropriate method to measure counterparty credit risk and
supplement it with stress testing. The aggregation of exposures should be carefully
reviewed to obtain a firm-wide view of risk.
\item[\textcolor{rust}{\sffamily\bfseries 10.}] \textbf{Interest rate risk in the
banking book (IRRBB).} Banks should carefully model embedded options in banking
products. They should also understand the trade-offs between earnings-based and
economic value-based approaches when measuring interest rate risk.
\end{enumerate}

% ---------------------------------------------------------------------
\lohead{58}{c}{Explain benefits and impacts of using an economic capital framework
within the following areas: credit portfolio management; risk-based pricing;
customer profitability analysis; and management incentives.}

\orsub{Benefits and impacts of using an economic capital framework}

\orsubb{1) Credit portfolio management}
\begin{checks}
\item A loan with a higher stand-alone risk does not necessarily contribute more
risk to the portfolio.
\item A loan's marginal contribution to the portfolio is therefore critical to
assessing the concentration of the portfolio.
\item Economic capital is a measurement of the level of concentration.
\end{checks}

\orsubb{2) Risk-based pricing}
\begin{checks}
\item Pricing decisions are based on expected risk-adjusted return on capital
(RAROC).
\item Pricing decisions include: (1) cost of funding, (2) expected loss,
(3) allocated economic capital, and (4) additional return required by
shareholders.
\item Can be used to maximise the bank's profitability.
\end{checks}

\orsubb{3) Customer profitability analysis}
\begin{checks}
\item The analysis is complicated in that many risks need to be aggregated at
customer level.
\item Customers need to be segmented in terms of ranges of net return per unit of
risk.
\item Economic capital is used in maximising the risk-return tradeoff.
\end{checks}

\orsubb{4) Management incentives}
\begin{checks}
\item Compensation schemes are a minor consideration in terms of the actual uses of
economic capital measures at the business unit level.
\item Management incentives is the issue that motivates bank managers to
participate in the technical aspects of the economic capital allocation process.
\end{checks}

% ---------------------------------------------------------------------
\lohead{58}{d}{Describe best practices and assess key concerns for the governance
of an economic capital framework.}

\orsub{Best practices for the governance of an economic capital framework}

\orsubb{1) Senior management commitment}
\begin{checks}
\item The successful implementation of an economic capital framework depends on the
involvement and experience of the senior management group.
\end{checks}

\orsubb{2) The business unit and its level of expertise}
\begin{checks}
\item Governance structures differ among banks.
\item Each business unit within the bank will manage its risk in accordance with
the amount of allocated capital.
\item The responsibilities for allocating capital within business units will also
vary among banks, as will the flexibility to reallocate capital during the
budgeting period.
\end{checks}

\orsubb{3) The timing of economic capital measurement and disclosures}
\begin{checks}
\item Most banks will compute economic capital on either a monthly or quarterly
basis.
\item Implementation of Basel II has fostered public disclosure of quantitative
information on economic capital measures among banks.
\end{checks}

\orsubb{4) Policies and procedures for owning, developing, validating, and monitoring economic capital models}
\begin{checks}
\item Formal policies and procedures encourage the consistent application of
economic capital across the bank.
\item The owner of the economic capital model will usually oversee the economic
capital framework.
\end{checks}

\orsub{Key concerns for governance and application}

\begin{enumerate}
\item[\textcolor{rust}{\sffamily\bfseries 1.}] \textbf{Senior management
commitment.} Economic capital is effective only when senior management fully
supports it and uses it for strategic planning and business decisions.
\item[\textcolor{rust}{\sffamily\bfseries 2.}] \textbf{Role of stress testing.}
Banks should use integrated stress tests to better assess the impact of adverse
scenarios on economic capital, rather than relying on isolated stress tests.
\item[\textcolor{rust}{\sffamily\bfseries 3.}] \textbf{Measuring risk.} Economic
capital can be measured on an absolute or relative basis. Its interpretation
depends on assumptions about diversification, management actions, and how well the
model captures risk.
\item[\textcolor{rust}{\sffamily\bfseries 4.}] \textbf{Economic capital is not the
only measure.} It should not be the sole determinant of required capital. Banks
must balance shareholders' preference for lower capital (higher profitability) with
rating agencies' preference for higher capital (greater solvency).
\item[\textcolor{rust}{\sffamily\bfseries 5.}] \textbf{Available capital
resources.} There is no standard definition of available capital across banks. Most
banks use adjusted Tier 1 capital as the basis for available capital.
\item[\textcolor{rust}{\sffamily\bfseries 6.}] \textbf{Transparency.} Economic
capital models should be transparent, well-documented, and easy to understand.
Greater transparency improves their credibility and usefulness in business
decision-making.
\end{enumerate}
